% SOURCE_AUTHORITY: sga5_fr_workpass.tex, lines 8707--8745 (LF-normalized unit).
% SOURCE_SHA256: 791F4EFFC5E02832D5D77ED03518C8156D6F07E4C8238B03545DB93D883FBB28
\subsubsection*{1.5. \((AR,\mathbb Z_\ell)\)-feixes construtíveis}

\begin{statement}{Definição 1.5.1}
Seja \(X\) um pré-esquema localmente noetheriano. Um objeto \(F\) de \(\Hom(\mathbb N^\circ,\Ab(X))\) chama-se \((AR,\mathbb Z_\ell)\)-feixe construtível se \(X\) admite um recobrimento
\[
(U_i)_{i\in I}
\]
Por abertos noetherianos tais que a restrição de \(F\) sobre cada \(U_i\) seja um objeto \((AR,\ell\mathbb Z)\)-noetheriano (V 1.3).
\end{statement}

Segue-se imediatamente de (V 5.2.1) que a subcategoria plena de \(\Hom(\mathbb N^\circ,\Ab(X))\) gerada pelos \((AR,\mathbb Z_\ell)\)-feixes construtíveis é estável por núcleos e conúcleos. Em particular, é uma categoria abeliana.

\subsubsection*{1.5.2. \((AR,\mathbb Z_\ell)\)-feixes construtíveis localmente \(AR\)-nulos}

É claro que um objeto \(AR\)-nulo (V 2.2.1) de \(\underline{\Hom}(\mathbb N^\circ,\Ab(X))\) é um \((AR,\mathbb Z_\ell)\)-feixe construtível se e somente se seus componentes são construtíveis. Segue-se, em particular, de (V 2.2.2) e (SGA 4 IX 2.6) que a categoria dos \((AR,\mathbb Z_\ell)\)-feixes construtíveis \(AR\)-nulos é uma subcategoria abeliana espessa da categoria plena engendrada pelos \((AR,\mathbb Z_\ell)\)-feixes construtíveis em \(\Hom(\mathbb N^\circ,\Ab(X))\).

Convindo em dizer que um objeto \(P\) de \(\underline{\Hom}(\mathbb N^\circ,\Ab(X))\) é localmente \(AR\)-nulo se existe um recobrimento aberto de \(X\) para a topologia étale (ou, o que equivale, para a topologia de Zariski) sobre cada um de cujos abertos \(P\) é \(AR\)-nulo, vê-se, reduzindo por localização ao caso anterior, que a categoria dos \((AR,\mathbb Z_\ell)\)-feixes construtíveis localmente \(AR\)-nulos é uma subcategoria abeliana espessa da dos \((AR,\mathbb Z_\ell)\)-feixes construtíveis.

Observe que, se \(X\) é noetheriano, os \((AR,\mathbb Z_\ell)\)-feixes construtíveis são os objetos \((AR,\ell\mathbb Z)\)-ádicos noetherianos de \(\Hom(\mathbb N^\circ,\Ab(X))\), e ser localmente \(AR\)-nulo equivale a ser \(AR\)-nulo.

\begin{statement}{Definição 1.5.3}
Seja \(X\) um pré-esquema localmente noetheriano. Chama-se categoria dos \((AR,\mathbb Z_\ell)\)-feixes construtíveis sobre \(X\), e denota-se por
\[
(AR,\mathbb Z_\ell)\text{-}\fc(X),
\]
a categoria abeliana quociente da subcategoria plena engendrada pelos \((AR,\mathbb Z_\ell)\)-feixes construtíveis em
\[
\underline{\Hom}(\mathbb N^\circ,\Ab(X))
\]
pela subcategoria plena, espessa, engendrada pelos \((AR,\mathbb Z_\ell)\)-feixes construtíveis localmente \(AR\)-nulos.
\end{statement}

\begin{statement}{Notação 1.5.4}
Sempre que for útil, denotar-se-á por
\[
U(X)
\]
a subcategoria plena de \(\Hom(\mathbb N^\circ,\Ab(X))\) engendrada pelos \((AR,\mathbb Z_\ell)\)-feixes construtíveis.
\end{statement}
